% Letter of Collaboration to an Institution (NON-BINDING). RFC-0136 D1, D7.
%
% Issued to an ORGANISATION COUNTERPARTY, not to a person (`recipient: organisation` in
% registry/lifecycle.yaml): `docket counterparty issue <slug> ICL`. The addressee block reads the
% counterparty's three facts -- \dkCOffice, \dkCLegalName, \dkCAddress -- and nothing about any
% person at the institution is held. The optional acknowledgement block is ruled lines, filled in
% by hand by whoever signs it.
%
% From knowledge docs/2026-11-07-personas-and-paridyum-in/08-paperwork/04, with the IIC/E-cell
% "body" placeholders folded into the receiving office and the institution's name, because a
% counterparty holds those three facts and no others.
%
% DRAFT FOR REVIEW BY QUALIFIED COUNSEL (06-counsel-brief.md Q11-Q13). Not legal advice. Do not
% issue to an institution before that review; HR-CHECKLIST.md states the order.
%
% WHY THESE TERMS. No money and no student data either way, no logo use without permission, no
% exclusivity, "being built" tense: BET-0010, RFC-0127 D4/D7, REQ-0003, STD-0021, RFC-0130 D7.
% Deliberately non-binding: PARIDYUM's CIN is null (registry/org.yaml), so it cannot yet hold the
% obligations a binding instrument creates. The binding form is the MOU, which waits for that.
\documentclass{paridyum}
\begin{document}

\paridyumletterhead
\paridyumdateref{\dkDocDate}{\dkDocRef}
\paridyumverifiable{\dkDocVerify}

\textbf{\dkCOffice} \\
\dkCLegalName \\
\dkCAddress

\vspace{0.9em}
\textbf{Subject: Letter of collaboration --- talks and workshops for student founders}

\vspace{0.7em}
Dear Sir or Madam,

\vspace{0.4em}
Thank you for speaking with us. This letter sets out what \textbf{\dkOrgName} would like to offer
the students of \dkCLegalName, and what we would ask in return. It is a statement of intent, not a
contract.

\section*{1. About Paridyum}
Paridyum is a technology venture building a platform on which a person can set up and run a
business. It is being built: we describe what exists, and we do not describe as available what
is not.

Paridyum is currently undergoing incorporation as a private limited company under the Companies
Act, 2013. Nothing in this letter is a representation that Paridyum has already been
incorporated as a separate legal entity.

\section*{2. What Paridyum offers}
\begin{itemize}[leftmargin=*, noitemsep]
    \item \textbf{Talks and workshops} for students who want to run a business --- for example,
          on choosing a legal form, handling money, GST thresholds, and setting up a business
          online --- at dates agreed between us. Our speakers describe Paridyum as it is today.
    \item \textbf{Written material} the institution may reuse with its students, such as slides
          or a one-page guide. It carries no tracking links.
    \item \textbf{The same terms as anyone else.} Students and the institution's own ventures may
          create organisations at paridyum.in on Paridyum's standard terms, in their own right.
          Nothing about the institution changes those terms, and no student is offered anything
          more, or less, because of this letter.
\end{itemize}

\section*{3. What we ask of the institution}
\begin{itemize}[leftmargin=*, noitemsep]
    \item A slot in the institution's calendar for each talk or workshop, and a room or an online
          meeting for it.
    \item A contact for arranging dates, by email to the address below.
\end{itemize}

\vspace{0.5em}
\begin{infobox}
\textbf{What neither side will do}\\[0.5em]
\begin{itemize}[leftmargin=*, noitemsep]
    \item \textbf{No money changes hands, in either direction.} Paridyum pays no sponsorship or
          fee, and charges none. Paridyum offers no prizes in credits, free months or capacity.
    \item \textbf{No student data is shared, ever, in either direction.} Paridyum will not ask
          for, accept, or keep any list of students, attendees, registrants or their contact
          details. Paridyum will not collect sign-ups at a session. A student who wants to use
          Paridyum goes to paridyum.in themselves.
    \item \textbf{No use of the other's name or logo without separate written permission.}
          Paridyum will not use the institution's name, logo or emblem, and the institution will
          not use Paridyum's, except as each side separately permits in writing.
    \item \textbf{No exclusivity.} Either side may work with anyone else.
\end{itemize}
\end{infobox}

\section*{4. Students under 18}
Paridyum's service is for people aged 18 and over. Sessions under this letter are for students
of the institution in its higher-education programmes. Paridyum will not run sessions at
schools under this letter.
% src: REQ-0038 item 1; RFC-0134 D9, D12. COUNSEL Q7.

\section*{5. Ending it}
Either side may end this collaboration at any time by email to the other, without giving a
reason. Paridyum's address for this is \dkOrgEmail. Nothing survives the end except the
commitments in the box above about student data and the use of names and logos.

\section*{6. Not a contract}
This letter records what each side intends. It does not create legal obligations, and neither
side may rely on it as a contract. If either side later wants binding terms, they will be in a
separate written agreement, which Paridyum can enter into only after its incorporation is
complete.

\vspace{0.7em}
If this is acceptable, please reply by email to \dkOrgEmail{} confirming it, with a contact for
scheduling. You may also confirm the authenticity of this letter using the reference and code
printed at the foot of this page.

\vspace{1.2em}
Yours faithfully,

\vspace{1.2em}
\paridyumsignblock

\vspace{1.5em}
\textbf{Acknowledged for \dkCLegalName} \textit{(optional; an email reply is enough)}

\vspace{0.5em}
\begin{tabularx}{\linewidth}{@{}F{3.4cm}X@{}}
Name:        & \dkblank[6cm] \\[1em]
Designation: & \dkblank[6cm] \\[1em]
Signature:   & \dkblank[6cm] \\[1em]
Date:        & \dkblank[6cm]
\end{tabularx}

\end{document}
